\section{Amplification and growth for one fixed divisor}
\label{sec:growth}

We now keep the entire divisor fixed while the support window grows.
Assume that \(Z\) is reflected, lies in \(|\Re z|\le a\), has
positive integer multiplicities, and satisfies
\begin{equation}\label{grow:count}
 N_0(Y)\le C(1+Y)^m\qquad(Y\ge0),
 \qquad C\ge1,\quad m\ge0.
\end{equation}
Suppose that its off-axis part is nonempty, and put
\begin{equation}\label{grow:definitions}
 s_Z=\sup\{\Re z:z\in Z,\ \Re z>0\},\qquad
 \cM_Z(L)=\sup_{0\ne f\in C_c^\infty((-L,L))}
       \left(-\frac{Q_Z(f,f)}{\norm f_2^2}\right)_+.
\end{equation}
Thus \(0<s_Z\le a\). The envelope is allowed to be infinite;
we interpret \(\log0=-\infty\) and \(\log(+\infty)=+\infty\).
The count hypothesis alone gives the lower rate
\begin{equation}\label{grow:lower-rate}
 \liminf_{L\to\infty}\frac{\log\cM_Z(L)}L\ge2s_Z.
\end{equation}
We first prove this assertion by retaining the free annihilation radius
in the support--margin construction. We then give a weighted local
condition under which the envelope is finite on every window and the
lower rate is exact. No arithmetic information is used in this section.

\subsection{Fixed-target amplification with arbitrarily small relative support loss}

First suppose that a target has been centered at the real pair
\(\pm\beta\), with \(0<\beta\le1\). Denote the centered divisor
by \(W\), and fix an annihilation radius \(R\ge16(a+2)\).
Apply the construction of \eqref{frontier:master} with
\(\ell=1/\beta\) and \(\tau=\beta T\), where \(T\) tends to
infinity. In particular \(\beta\ell=1\), and \(\tau\ge1\)
once \(T\ge1/\beta\). All counting budgets and all constants
attached to this fixed \(R\) are fixed before \(T\) increases.
The exponent estimate \eqref{frontier:K} has the form
\begin{equation}\label{grow:fixed-radius-K}
 K\le a(T+\ell)+(M-1)
       \{\log(1/\beta)+\log(1+\beta T)\}+J_R,
\end{equation}
where \(M\) is a fixed finite low-node budget and \(J_R\) is
independent of \(T\). The exact support estimate
\eqref{frontier:support-exact} consequently gives
\begin{equation}\label{grow:fixed-radius-support}
 L_T\le\ell+T+4K/R
 \le (1+4a/R)T+A_R\log(2+T)+B_R.
\end{equation}
Here and below the finite constants may depend on the fixed target,
the strip, the count, and \(R\). This is not a uniform estimate
over target heights or over growing values of \(R\).

We spell out why the same construction gives a normalized, rather
than merely unnormalized, exponential gain. Its local target values
before translation are \(\sqrt\ell\) and \(-\sqrt\ell\).
Multiplication by \(\cosh(Tw)\) is the average of translations by
\(T\) and \(-T\), so it does not increase the \(L^2\) norm.
The chosen interpolation components have imaginary width
\(O((\ell+T)^{-1})\). Thus every selected reflected pair retains
a negative contribution after translation, while all intermediate
nodes vanish. The absolute tail is at most \(\ell/2\). In
particular the actual target, of multiplicity at least one, yields
\begin{equation}\label{grow:fixed-radius-raw}
 Q_W(f_T,f_T)\le-2\ell\cosh^2(\beta T)+\ell/2
                 \le-\frac\ell4 e^{2\beta T}.
\end{equation}
These are the sign and tail conclusions of
\eqref{frontier:master}; they do not follow from translating an
arbitrary previously negative test.

The physical-coordinate norm estimate in that construction is
\begin{equation}\label{grow:fixed-radius-norm}
 \norm{f_T}_2^2
 \le D_R(1+\beta T)^{2M-2}
             \exp\left(\frac{D_R\beta^2K}{(1+\beta T)^2}\right).
\end{equation}
The constant in the exponent can be enlarged to the same symbol
\(D_R\). Equation \eqref{grow:fixed-radius-K} shows that this
exponent remains bounded as \(T\to\infty\), and in fact tends
to zero for a fixed target. Combining
\eqref{grow:fixed-radius-raw} and \eqref{grow:fixed-radius-norm}
therefore gives a nonzero compact smooth test for each sufficiently
large \(T\), with
\begin{equation}\label{grow:fixed-radius-margin}
 -\frac{Q_W(f_T,f_T)}{\norm{f_T}_2^2}
       \ge c_R\frac{e^{2\beta T}}{(1+T)^{2M-2}},
 \qquad c_R>0.
\end{equation}
Its support is strictly inside the interval with radius bounded in
\eqref{grow:fixed-radius-support}, because the initial smoothing
function has support strictly inside its prescribed interval.

Fix \(0<\delta<1\) and, for each requested radius \(L\), set
\[
 T=\frac{(1-\delta)L}{1+4a/R}.
\]
For every sufficiently large \(L\), the logarithmic error in
\eqref{grow:fixed-radius-support} is less than \(\delta L\),
and \(T\ge1/\beta\). The constructed test then lies in
\(C_c^\infty((-L,L))\). Equation
\eqref{grow:fixed-radius-margin} implies
\[
 \liminf_{L\to\infty}\frac{\log\cM_W(L)}L
       \ge\frac{2\beta(1-\delta)}{1+4a/R}.
\]
First let \(\delta\downarrow0\), and then let the fixed radius
\(R\) increase without bound. This proves a lower rate
\(2\beta\). The order of limits matters: no control of
\(A_R,B_R,c_R,D_R,J_R\) as \(R\to\infty\) has been asserted
or used. The construction supplies tests on every sufficiently
large window, not just a subsequence.

For a target \(\beta+i\gamma\) in \(Z\), centering changes the
count only to
\begin{equation}\label{grow:shifted-count}
 N_\gamma(Y)\le C(1+|\gamma|+Y)^m
       \le C(1+|\gamma|)^m(1+Y)^m.
\end{equation}
If \(W=Z-i\gamma\) and \(f(u)=e^{-i\gamma u}g(u)\), then
\(H_f(z)=H_g(z-i\gamma)\), so
\(Q_Z(f,f)=Q_W(g,g)\), with the same norm and support. Hence
the preceding conclusion applies at any fixed height.

The restriction \(\beta\le1\) is immaterial for the growth
statement. For \(\alpha=\max\{a,1\}\), let
\(\widetilde Z=\{z/\alpha:z\in Z\}\), retaining multiplicities.
Its strip width is at most one and its count is bounded by
\(C\alpha^m(1+Y)^m\). If
\(f(u)=\sqrt\alpha\,g(\alpha u)\), direct substitution gives
\begin{equation}\label{grow:dilation}
 \norm f_2=\norm g_2,\qquad
 H_f(z)=\alpha^{-1/2}H_g(z/\alpha),\qquad
 Q_Z(f,f)=\alpha^{-1}Q_{\widetilde Z}(g,g).
\end{equation}
A radius \(\alpha L\) for \(g\) becomes radius \(L\) for
\(f\); the inverse scaling gives the exact envelope relation
\(\cM_Z(L)=\alpha^{-1}\cM_{\widetilde Z}(\alpha L)\).
Thus the lower rate for a target of displacement
\(\beta/\alpha\) scales back to \(2\beta\).

We have proved the lower rate for every fixed off-axis target.
Taking their supremum proves \eqref{grow:lower-rate}, even when
\(s_Z\) is not attained. In particular \(\cM_Z(L)>0\) for
every sufficiently large \(L\), and it tends to infinity.
For any finite requested margin, a supremum strictly greater than
that margin contains an actual compact smooth test above it. Thus
all finite margins are attained as inequalities on some finite
window, without an eigenfunction-attainment assumption.

\subsection{A weighted local condition sufficient for a finite envelope}

Use one positive-real representative of each reflected pair, and assume
\begin{equation}\label{grow:weighted-density}
 B_Z:=\sup_{y\in\R}
 \sum_{\substack{z=\sigma+i\gamma\in Z\\
                  \sigma>0,\ |\gamma-y|\le1}}
          \nu(z)\sigma^2<\infty.
\end{equation}
This condition permits infinitely many off-axis pairs, unbounded
unweighted local counts when the displacements shrink, and arbitrarily
close distinct locations. It holds for a finite off-axis part or a
uniformly bounded unit-height count. The axis part requires no extra
hypothesis beyond \eqref{grow:count}. We will prove, for every
\(L>0\) and \(0<\varepsilon\le1\),
\begin{equation}\label{grow:epsilon-upper}
 \cM_Z(L)\le
 8B_Z(s_Z+\varepsilon)\varepsilon^{-2}L^2
                   e^{2(s_Z+\varepsilon)L}.
\end{equation}
In particular, for \(L\ge1\),
\begin{equation}\label{grow:upper}
 \cM_Z(L)\le8e^2B_Z(s_Z+1)L^4e^{2s_ZL}.
\end{equation}
No optimality of the polynomial factor is asserted.

Fix a compact smooth test \(f\) supported in \((-L,L)\), and set
\[
 A_z=\frac{H_f(\sigma+i\gamma)+H_f(-\sigma+i\gamma)}2,
 \qquad
 D_z=\frac{H_f(\sigma+i\gamma)-H_f(-\sigma+i\gamma)}2.
\]
The contribution of the pair is exactly
\(2\nu(z)(|A_z|^2-|D_z|^2)\); every axis term is
nonnegative. The entire function
\[
 G_f(w)=\int_\R u f(u)e^{wu}\,du
\]
satisfies
\[
 D_z=\frac\sigma2\int_0^1
     \{G_f(t\sigma+i\gamma)+G_f(-t\sigma+i\gamma)\}\,dt.
\]
Cauchy--Schwarz on this average gives
\begin{equation}\label{grow:pair-bound}
 2\nu(z)|D_z|^2\le\nu(z)\sigma^2\int_0^1
       \{|G_f(t\sigma+i\gamma)|^2+
          |G_f(-t\sigma+i\gamma)|^2\}\,dt.
\end{equation}

For fixed \(t\in[0,1]\), give each of \(t\sigma+i\gamma\)
and \(-t\sigma+i\gamma\) the weight \(\nu(z)\sigma^2\), and
denote the resulting measure by \(\mu_t\). At \(t=0\), both
weights are retained even though their locations coincide. Its mass
in any radius-one vertical band is at most \(2B_Z\). The area
mean inequality for the subharmonic function \(|G|^2\) yields
\begin{equation}\label{grow:area-mean}
 \begin{split}
 \int |G(w)|^2\,d\mu_t(w)
 &\le\frac1{\pi\varepsilon^2}
    \int\!\!\int_{|\zeta-w|<\varepsilon}
               |G(\zeta)|^2\,dA(\zeta)\,d\mu_t(w)\\
 &\le\frac{2B_Z}{\pi\varepsilon^2}
    \int_{-s_Z-\varepsilon}^{s_Z+\varepsilon}
    \int_\R |G(x+iy)|^2\,dy\,dx.
 \end{split}
\end{equation}
Indeed, for a fixed area point \(x+iy\), all contributing centers
have \(|\gamma-y|<\varepsilon\le1\); their total weight is at
most \(2B_Z\). Their real parts lie in \([-s_Z,s_Z]\).
Tonelli justifies the interchange of these nonnegative integrands,
and no separation assumption enters the argument.

For \(G=G_f\), Plancherel gives
\[
 \int_\R|G_f(x+iy)|^2\,dy
   =2\pi\int_\R u^2|f(u)|^2e^{2xu}\,du
   \le2\pi L^2e^{2|x|L}\norm f_2^2.
\]
Discard the positive terms of \(Q_Z\), sum
\eqref{grow:pair-bound}, and apply \eqref{grow:area-mean}.
The real integral has length \(2(s_Z+\varepsilon)\), and the
\(t\)-integral has length one. This proves
\eqref{grow:epsilon-upper} with its displayed constant.
Taking \(\varepsilon=1/L\) proves \eqref{grow:upper}; taking
\(\varepsilon=1\) proves finiteness also for smaller windows.
The original form converges absolutely by \eqref{grow:count}.
Only its nonnegative majorants were interchanged in the proof.
There is no need to assert a bounded positive axis operator on the
whole \(L^2\) window.

\subsection{Exact exponential rate and the first prescribed-margin radius}

Combining \eqref{grow:lower-rate} and \eqref{grow:upper} gives
\begin{equation}\label{grow:exact-rate}
 \lim_{L\to\infty}\frac{\log\cM_Z(L)}L=2s_Z
 \qquad\text{if }B_Z<\infty.
\end{equation}
The logarithm is finite and ordinary for all sufficiently large
\(L\): finiteness follows from the weighted bound and positivity
from the target construction. Define the first radius using actual
tests by
\begin{equation}\label{grow:first-definition}
 r_Z(\eta)=\inf\{L>0:\ \exists\,0\ne f\in
 C_c^\infty((-L,L)),\quad Q_Z(f,f)\le-\eta\norm f_2^2\},
 \qquad\eta>0.
\end{equation}
It is finite, and
\begin{equation}\label{grow:first-asymptotic}
 r_Z(\eta)\sim\frac{\log\eta}{2s_Z}
                    \qquad(\eta\to\infty).
\end{equation}
Here the divisor, including all its multiplicities, is fixed.

To justify inversion without attainment, fix \(0<\delta<2s_Z\).
For all sufficiently large \(L\), the rate law implies
\[
 e^{(2s_Z-\delta)L}\le\cM_Z(L)
                        \le e^{(2s_Z+\delta)L}.
\]
At a radius strictly below \(\log\eta/(2s_Z+\delta)\), the
upper bound excludes a test of margin \(\eta\). Bounded smaller
radii are excluded for large \(\eta\) by monotonicity and
finiteness of the envelope. At a radius strictly above
\(\log\eta/(2s_Z-\delta)\), the lower bound makes the supremum
strictly exceed \(\eta\), hence gives an actual test with
strictly larger margin. Let the bounding radii approach their
endpoints and then let \(\delta\downarrow0\). This proves
\eqref{grow:first-asymptotic}, also when a strict margin inequality
is used in \eqref{grow:first-definition}.

The weighted condition is sufficient for the matching upper rate;
it is not asserted to be necessary for that rate. It does not
characterize all locally semibounded divisors.
The lower rate \eqref{grow:lower-rate} needs only the polynomial
count, whereas a finite upper rate needs further control. In
particular, local semiboundedness alone cannot be substituted for
\eqref{grow:weighted-density}; the screened construction in
Section~\ref{sec:arbitrary-growth} shows that fixed sparse counts can coexist with
arbitrarily rapid growth of the finite negative envelope.
